% !TEX program = xelatex
\def\StudyRoot{../..}
\def\StudyBook{Compact}
\documentclass[11pt,lang=cn,chinesefont=nofont,mode=simple,math=newtx,color=blue]{elegantbook}
\usepackage{longtable,array,fvextra,xurl,adjustbox,needspace}
\usepackage[normalem]{ulem}
\providecommand{\lt}{<}
\providecommand{\gt}{>}
\tracinglostchars=3
\setCJKmainfont{FandolSong-Regular.otf}[BoldFont=FandolSong-Bold.otf,ItalicFont=FandolKai-Regular.otf]
\setCJKsansfont{FandolHei-Regular.otf}[BoldFont=FandolHei-Bold.otf]
\setCJKmonofont{FandolFang-Regular.otf}
\setmonofont{DejaVuSansMono.ttf}[BoldFont=DejaVuSansMono-Bold.ttf,ItalicFont=DejaVuSansMono-Oblique.ttf,Scale=.84]
\geometry{a4paper,left=22mm,right=18mm,top=20mm,bottom=21mm,headheight=15pt,headsep=7mm}
\definecolor{structurecolor}{HTML}{173F57}
\colorlet{main}{structurecolor}
\colorlet{second}{structurecolor}
\colorlet{third}{structurecolor}
\hypersetup{unicode=true,colorlinks=true,linkcolor=structurecolor,urlcolor=structurecolor,citecolor=structurecolor,pdfauthor={CityU StudyNotes}}
\makeatletter
\renewcommand{\tableofcontents}{%
  \hypersetup{linktoc=all,linkcolor=black}%
  \chapter*{\contentsname}%
  \StudyTarget{\StudyBook__contents}\StudyTarget{\StudyBook___1}%
  \@starttoc{toc}%
  \hypersetup{linkcolor=structurecolor}\clearpage}
\makeatother
\setstretch{1.25}
\setlength{\parskip}{.35\baselineskip}
\setlength{\parindent}{0pt}
\setlength{\emergencystretch}{2em}
\setcounter{tocdepth}{1}
\setcounter{secnumdepth}{-1}
\setlist{nosep,leftmargin=1.8em,topsep=.35\baselineskip}
\titleformat{\chapter}[display]{\Huge\sffamily\bfseries\color{structurecolor}}{}{0pt}{}[\vspace{.4ex}\titlerule]
\titlespacing*{\chapter}{0pt}{0pt}{22pt}
\titleformat{\section}{\Large\sffamily\bfseries\color{structurecolor}}{}{0pt}{}
\titleformat{\subsection}{\large\sffamily\bfseries}{}{0pt}{}
\titleformat{\subsubsection}{\normalsize\sffamily\bfseries}{}{0pt}{}
\titlespacing*{\section}{0pt}{1.6\baselineskip}{.5\baselineskip}
\titlespacing*{\subsection}{0pt}{1.1\baselineskip}{.35\baselineskip}
\providecommand{\StudyBook}{Sample}
\providecommand{\StudyRoot}{../..}
\providecommand{\StudySourceRef}{main}
\providecommand{\StudyDate}{2026-10-16}
\newcommand{\StudyRunning}{\StudyBook}
\fancyhf{}
\fancyhead[L]{\small\sffamily\color{gray}\StudyRunning}
\fancyhead[R]{\small\sffamily\color{gray}CityU StudyNotes}
\fancyfoot[R]{\small\thepage}
\renewcommand{\headrulewidth}{.3pt}
\pagestyle{fancy}
\fancypagestyle{plain}{\fancyhf{}\fancyfoot[R]{\small\thepage}\renewcommand{\headrulewidth}{0pt}}
\newwrite\StudyMap
\AtBeginDocument{\immediate\openout\StudyMap=\jobname.pages}
\makeatletter
\newcommand{\StudyTarget}[1]{%
  \hypertarget{#1}{}\label{#1}%
  \protected@write\StudyMap{}{#1|\thepage}%
}
\makeatother
\newcommand{\StudyChapter}[2]{%
  \clearpage\renewcommand{\StudyRunning}{\ifcsname studytitle@#1\endcsname\csname studytitle@#1\endcsname\else#1\fi}%
  \chapter*{#2}\phantomsection\addcontentsline{toc}{chapter}{#2}%
  \StudyTarget{#1__root}%
}
\newcommand{\StudyHeading}[2]{%
  \ifcase#1\relax\or
    \section*{#2}\phantomsection\addcontentsline{toc}{section}{#2}%
  \or\subsection*{#2}\or\subsubsection*{#2}\else\paragraph*{#2}\fi
}
\newcommand{\StudyImage}[1]{%
  \par\medskip\begingroup
  \edef\StudyResolvedImage{\ifcsname studyimage@#1\endcsname\csname studyimage@#1\endcsname\else\StudyRoot/#1\fi}%
  \noindent\makebox[\linewidth][c]{%
  \adjustbox{max width=.96\linewidth,max height=.69\textheight}{%
  \includegraphics{\StudyResolvedImage}}}\endgroup\par\medskip
}
\newcommand{\StudyCellImage}[1]{%
  \begingroup
  \edef\StudyResolvedImage{\ifcsname studyimage@#1\endcsname\csname studyimage@#1\endcsname\else\StudyRoot/#1\fi}%
  \adjustbox{max width=\linewidth,max height=.28\textheight}{\includegraphics{\StudyResolvedImage}}\endgroup}
\DefineVerbatimEnvironment{StudyCode}{Verbatim}{breaklines=true,breakanywhere=true,fontsize=\small,tabsize=4,framesep=3mm,frame=leftline,rulecolor=\color{gray!55}}
\newcommand{\StudyInline}[1]{\allowbreak{}\texttt{#1}\allowbreak{}}
\newenvironment{StudySolution}[1]{%
  \par\medskip\Needspace{3\baselineskip}\noindent
  {\bfseries\color{structurecolor}#1}\par\nobreak\smallskip
}{\par\medskip}
\newcommand{\StudyBadge}[2]{%
  \begingroup\small\sffamily
  \ifstrequal{#1}{mid}{\colorlet{studybadge}{green!35!black}}{%
  \ifstrequal{#1}{final}{\colorlet{studybadge}{yellow!45!black}}{\colorlet{studybadge}{red!45!black}}}%
  \textcolor{studybadge}{\rule[-.2ex]{1.2pt}{1.2em}}\enspace\textcolor{black}{#2}\quad\endgroup}
\newcommand{\StudyRef}[3]{%
  \ifcsname studyref@#1@#2\endcsname
    \csname studyref@#1@#2\endcsname{#3}%
  \else #3\textsuperscript{[待定页]}\fi}
\newcommand{\StudyRepoRef}[3]{%
  \href{https://github.com/CrazyShout/CityU-StudyNotes/blob/\StudySourceRef/#1\ifstrempty{#2}{}{\##2}}{#3}}
\newcommand{\StudyNotebookCell}[2]{%
  \ifcsname studynb@#1@#2\endcsname
    \input{\csname studynb@#1@#2\endcsname}%
  \else\PackageError{studynotes}{Missing saved notebook cell #2}{}\fi}
\newcommand{\StudyCoverExtra}{}
\newcommand{\StudyCover}[2]{%
  \thispagestyle{empty}
  \StudyTarget{\StudyBook__root}\StudyCoverExtra
  \vspace*{22mm}
  {\sffamily\large\color{structurecolor}CITYU STUDYNOTES\par}
  \vspace{12mm}
  {\sffamily\bfseries\Huge\color{structurecolor}#1\par}
  \vspace{6mm}{\Large #2\par}
  \vspace{10mm}\noindent\textcolor{structurecolor}{\rule{.42\linewidth}{1pt}}\par
  \vspace{8mm}中文讲解 · English terminology · 双语题答\par
  A4 单栏 · 根据已下载教师材料整理\par
  \vfill
  {\small 本资料用于课程学习与复习，保留原材料、历史题源与教学补充的身份。\par
  模板：ElegantBook 4.7；版式与构建由 CityU-StudyNotes 维护。\par
  \StudyDate\par}
  \clearpage
}
\newcommand{\StudyQuestionExtra}{}
\newcommand{\StudyQuestion}[3]{%
  \Needspace{10\baselineskip}
  \section*{#1 · #2}\phantomsection\addcontentsline{toc}{section}{#1 · #2}%
  \StudyTarget{\StudyBook__#1}%
  \StudyQuestionExtra
  \lowercase{\def\StudyQuestionAnchor{#1}}%
  \noindent\begin{minipage}[t]{.65\linewidth}{\small\sffamily #3}\end{minipage}\hfill
  \begin{minipage}[t]{.33\linewidth}\raggedleft\scriptsize
  \StudyRef{Questions}{\StudyQuestionAnchor}{题目} · \StudyRef{Answers}{\StudyQuestionAnchor}{答案}
  \end{minipage}\par\medskip
}
\InputIfFileExists{\StudyRoot/.build/latex/\StudyBook/refs.tex}{}{}

\begin{document}
\setlength{\parskip}{.25\baselineskip}
\titlespacing*{\subsection}{0pt}{.7\baselineskip}{.2\baselineskip}
\renewcommand{\StudyCoverExtra}{\StudyTarget{Compact__cs5489-concise-bilingual-questions-and-answers}}
\StudyCover{CS5489 · 期中精简题答}{Concise Bilingual Questions and Answers}
\input{\StudyRoot/CS5489/exam-review/Midterm/frontmatter/Compact.tex}
\begingroup\small\setlength{\cftbeforechapskip}{.45\baselineskip}\tableofcontents\endgroup
\Needspace{10\baselineskip}\phantomsection\addcontentsline{toc}{chapter}{学习流程与模型诊断}
\renewcommand{\StudyQuestionExtra}{\StudyTarget{Compact__mt011}}
\StudyQuestion{MT011}{识别过拟合，并提出可操作的改进}{简答题 / Written response · 解释、提出办法}
\input{\StudyRoot/CS5489/exam-review/Midterm/items/MT011/provenance.tex}
\subsection*{English question}
\StudyTarget{Compact__english-question}
\input{\StudyRoot/CS5489/exam-review/Midterm/items/MT011/question-en.tex}
\subsection*{中文题意}
\StudyTarget{Compact___2}
\input{\StudyRoot/CS5489/exam-review/Midterm/items/MT011/question-zh.tex}
\subsection*{English answer}
\StudyTarget{Compact__english-answer}
\input{\StudyRoot/CS5489/exam-review/Midterm/items/MT011/answer-en.tex}
\subsection*{中文对应答案}
\StudyTarget{Compact___3}
\input{\StudyRoot/CS5489/exam-review/Midterm/items/MT011/answer-zh.tex}
\input{\StudyRoot/CS5489/exam-review/Midterm/items/MT011/concise-extra.tex}
\renewcommand{\StudyQuestionExtra}{\StudyTarget{Compact__mt004}\StudyTarget{Compact__mt004-cv}}
\StudyQuestion{MT004}{CV选完参数后，为什么还要重新训练？}{选择题（多选） / Multiple select}
\input{\StudyRoot/CS5489/exam-review/Midterm/items/MT004/provenance.tex}
\subsection*{English question}
\StudyTarget{Compact__english-question_1}
\input{\StudyRoot/CS5489/exam-review/Midterm/items/MT004/question-en.tex}
\subsection*{中文题意}
\StudyTarget{Compact___4}
\input{\StudyRoot/CS5489/exam-review/Midterm/items/MT004/question-zh.tex}
\subsection*{English answer}
\StudyTarget{Compact__english-answer_1}
\input{\StudyRoot/CS5489/exam-review/Midterm/items/MT004/answer-en.tex}
\subsection*{中文对应答案}
\StudyTarget{Compact___5}
\input{\StudyRoot/CS5489/exam-review/Midterm/items/MT004/answer-zh.tex}
\input{\StudyRoot/CS5489/exam-review/Midterm/items/MT004/concise-extra.tex}
\renewcommand{\StudyQuestionExtra}{\StudyTarget{Compact__mt028}}
\StudyQuestion{MT028}{训练误差低、测试误差高：可能原因}{选择题（多选） / Multiple select}
\input{\StudyRoot/CS5489/exam-review/Midterm/items/MT028/provenance.tex}
\subsection*{English question}
\StudyTarget{Compact__english-question_2}
\input{\StudyRoot/CS5489/exam-review/Midterm/items/MT028/question-en.tex}
\subsection*{中文题意}
\StudyTarget{Compact___6}
\input{\StudyRoot/CS5489/exam-review/Midterm/items/MT028/question-zh.tex}
\subsection*{English answer}
\StudyTarget{Compact__english-answer_2}
\input{\StudyRoot/CS5489/exam-review/Midterm/items/MT028/answer-en.tex}
\subsection*{中文对应答案}
\StudyTarget{Compact___7}
\input{\StudyRoot/CS5489/exam-review/Midterm/items/MT028/answer-zh.tex}
\input{\StudyRoot/CS5489/exam-review/Midterm/items/MT028/concise-extra.tex}
\renewcommand{\StudyQuestionExtra}{\StudyTarget{Compact__mt002}}
\StudyQuestion{MT002}{训练准、测试差：哪些改进值得尝试？}{选择题（多选） / Multiple select}
\input{\StudyRoot/CS5489/exam-review/Midterm/items/MT002/provenance.tex}
\subsection*{English question}
\StudyTarget{Compact__english-question_3}
\input{\StudyRoot/CS5489/exam-review/Midterm/items/MT002/question-en.tex}
\subsection*{中文题意}
\StudyTarget{Compact___8}
\input{\StudyRoot/CS5489/exam-review/Midterm/items/MT002/question-zh.tex}
\subsection*{English answer}
\StudyTarget{Compact__english-answer_3}
\input{\StudyRoot/CS5489/exam-review/Midterm/items/MT002/answer-en.tex}
\subsection*{中文对应答案}
\StudyTarget{Compact___9}
\input{\StudyRoot/CS5489/exam-review/Midterm/items/MT002/answer-zh.tex}
\input{\StudyRoot/CS5489/exam-review/Midterm/items/MT002/concise-extra.tex}
\renewcommand{\StudyQuestionExtra}{\StudyTarget{Compact__mt020}}
\StudyQuestion{MT020}{数据很多，训练和测试仍都差}{简答题 / Written response · 诊断问题、提出办法}
\input{\StudyRoot/CS5489/exam-review/Midterm/items/MT020/provenance.tex}
\subsection*{English question}
\StudyTarget{Compact__english-question_4}
\input{\StudyRoot/CS5489/exam-review/Midterm/items/MT020/question-en.tex}
\subsection*{中文题意}
\StudyTarget{Compact___10}
\input{\StudyRoot/CS5489/exam-review/Midterm/items/MT020/question-zh.tex}
\subsection*{English answer}
\StudyTarget{Compact__english-answer_4}
\input{\StudyRoot/CS5489/exam-review/Midterm/items/MT020/answer-en.tex}
\subsection*{中文对应答案}
\StudyTarget{Compact___11}
\input{\StudyRoot/CS5489/exam-review/Midterm/items/MT020/answer-zh.tex}
\input{\StudyRoot/CS5489/exam-review/Midterm/items/MT020/concise-extra.tex}
\renewcommand{\StudyQuestionExtra}{\StudyTarget{Compact__mt022}}
\StudyQuestion{MT022}{五万客户、需要解释：选择哪种分类器？}{简答题 / Written response · 比较、提出办法}
\input{\StudyRoot/CS5489/exam-review/Midterm/items/MT022/provenance.tex}
\subsection*{English question}
\StudyTarget{Compact__english-question_5}
\input{\StudyRoot/CS5489/exam-review/Midterm/items/MT022/question-en.tex}
\subsection*{中文题意}
\StudyTarget{Compact___12}
\input{\StudyRoot/CS5489/exam-review/Midterm/items/MT022/question-zh.tex}
\subsection*{English answer}
\StudyTarget{Compact__english-answer_5}
\input{\StudyRoot/CS5489/exam-review/Midterm/items/MT022/answer-en.tex}
\subsection*{中文对应答案}
\StudyTarget{Compact___13}
\input{\StudyRoot/CS5489/exam-review/Midterm/items/MT022/answer-zh.tex}
\input{\StudyRoot/CS5489/exam-review/Midterm/items/MT022/concise-extra.tex}
\Needspace{10\baselineskip}\phantomsection\addcontentsline{toc}{chapter}{Bayes、Naive Bayes与模型比较}
\renewcommand{\StudyQuestionExtra}{\StudyTarget{Compact__mt040}\StudyTarget{Compact__l2q05}\StudyTarget{Compact__mt040-bayesnaive-bayes}}
\StudyQuestion{MT040}{Bayes规则和Naive Bayes的关系}{选择题（多选） / Multiple select}
\input{\StudyRoot/CS5489/exam-review/Midterm/items/MT040/provenance.tex}
\subsection*{English question}
\StudyTarget{Compact__english-question_6}
\input{\StudyRoot/CS5489/exam-review/Midterm/items/MT040/question-en.tex}
\subsection*{中文题意}
\StudyTarget{Compact___14}
\input{\StudyRoot/CS5489/exam-review/Midterm/items/MT040/question-zh.tex}
\subsection*{English answer}
\StudyTarget{Compact__english-answer_6}
\input{\StudyRoot/CS5489/exam-review/Midterm/items/MT040/answer-en.tex}
\subsection*{中文对应答案}
\StudyTarget{Compact___15}
\input{\StudyRoot/CS5489/exam-review/Midterm/items/MT040/answer-zh.tex}
\input{\StudyRoot/CS5489/exam-review/Midterm/items/MT040/concise-extra.tex}
\renewcommand{\StudyQuestionExtra}{\StudyTarget{Compact__mt052}\StudyTarget{Compact__l2q01}}
\StudyQuestion{MT052}{生成式分类器学了什么}{选择题（多选） / Multiple select}
\input{\StudyRoot/CS5489/exam-review/Midterm/items/MT052/provenance.tex}
\subsection*{English question}
\StudyTarget{Compact__english-question_7}
\input{\StudyRoot/CS5489/exam-review/Midterm/items/MT052/question-en.tex}
\subsection*{中文题意}
\StudyTarget{Compact___16}
\input{\StudyRoot/CS5489/exam-review/Midterm/items/MT052/question-zh.tex}
\subsection*{English answer}
\StudyTarget{Compact__english-answer_7}
\input{\StudyRoot/CS5489/exam-review/Midterm/items/MT052/answer-en.tex}
\subsection*{中文对应答案}
\StudyTarget{Compact___17}
\input{\StudyRoot/CS5489/exam-review/Midterm/items/MT052/answer-zh.tex}
\input{\StudyRoot/CS5489/exam-review/Midterm/items/MT052/concise-extra.tex}
\renewcommand{\StudyQuestionExtra}{\StudyTarget{Compact__mt046}\StudyTarget{Compact__l2q09}\StudyTarget{Compact__mt046-gaussian-nb}}
\StudyQuestion{MT046}{完整讲出Gaussian NB}{简答题 / Written response · 解释}
\input{\StudyRoot/CS5489/exam-review/Midterm/items/MT046/provenance.tex}
\subsection*{English question}
\StudyTarget{Compact__english-question_8}
\input{\StudyRoot/CS5489/exam-review/Midterm/items/MT046/question-en.tex}
\subsection*{中文题意}
\StudyTarget{Compact___18}
\input{\StudyRoot/CS5489/exam-review/Midterm/items/MT046/question-zh.tex}
\subsection*{English answer}
\StudyTarget{Compact__english-answer_8}
\input{\StudyRoot/CS5489/exam-review/Midterm/items/MT046/answer-en.tex}
\subsection*{中文对应答案}
\StudyTarget{Compact___19}
\input{\StudyRoot/CS5489/exam-review/Midterm/items/MT046/answer-zh.tex}
\input{\StudyRoot/CS5489/exam-review/Midterm/items/MT046/concise-extra.tex}
\renewcommand{\StudyQuestionExtra}{\StudyTarget{Compact__mt015}\StudyTarget{Compact__l2q08}\StudyTarget{Compact__mt015-nb}}
\StudyQuestion{MT015}{NB的最优性、独立性与边界}{选择题（多选） / Multiple select}
\input{\StudyRoot/CS5489/exam-review/Midterm/items/MT015/provenance.tex}
\subsection*{English question}
\StudyTarget{Compact__english-question_9}
\input{\StudyRoot/CS5489/exam-review/Midterm/items/MT015/question-en.tex}
\subsection*{中文题意}
\StudyTarget{Compact___20}
\input{\StudyRoot/CS5489/exam-review/Midterm/items/MT015/question-zh.tex}
\subsection*{English answer}
\StudyTarget{Compact__english-answer_9}
\input{\StudyRoot/CS5489/exam-review/Midterm/items/MT015/answer-en.tex}
\subsection*{中文对应答案}
\StudyTarget{Compact___21}
\input{\StudyRoot/CS5489/exam-review/Midterm/items/MT015/answer-zh.tex}
\input{\StudyRoot/CS5489/exam-review/Midterm/items/MT015/concise-extra.tex}
\renewcommand{\StudyQuestionExtra}{\StudyTarget{Compact__mt032}\StudyTarget{Compact__l2q06}\StudyTarget{Compact__mt032-bayesian-classifier}}
\StudyQuestion{MT032}{Bayesian classifier的性质}{选择题（多选） / Multiple select}
\input{\StudyRoot/CS5489/exam-review/Midterm/items/MT032/provenance.tex}
\subsection*{English question}
\StudyTarget{Compact__english-question_10}
\input{\StudyRoot/CS5489/exam-review/Midterm/items/MT032/question-en.tex}
\subsection*{中文题意}
\StudyTarget{Compact___22}
\input{\StudyRoot/CS5489/exam-review/Midterm/items/MT032/question-zh.tex}
\subsection*{English answer}
\StudyTarget{Compact__english-answer_10}
\input{\StudyRoot/CS5489/exam-review/Midterm/items/MT032/answer-en.tex}
\subsection*{中文对应答案}
\StudyTarget{Compact___23}
\input{\StudyRoot/CS5489/exam-review/Midterm/items/MT032/answer-zh.tex}
\input{\StudyRoot/CS5489/exam-review/Midterm/items/MT032/concise-extra.tex}
\renewcommand{\StudyQuestionExtra}{\StudyTarget{Compact__mt005}\StudyTarget{Compact__l2q07}\StudyTarget{Compact__mt005-nb}}
\StudyQuestion{MT005}{NB的学习、不确定性和正则化}{选择题（多选） / Multiple select}
\input{\StudyRoot/CS5489/exam-review/Midterm/items/MT005/provenance.tex}
\subsection*{English question}
\StudyTarget{Compact__english-question_11}
\input{\StudyRoot/CS5489/exam-review/Midterm/items/MT005/question-en.tex}
\subsection*{中文题意}
\StudyTarget{Compact___24}
\input{\StudyRoot/CS5489/exam-review/Midterm/items/MT005/question-zh.tex}
\subsection*{English answer}
\StudyTarget{Compact__english-answer_11}
\input{\StudyRoot/CS5489/exam-review/Midterm/items/MT005/answer-en.tex}
\subsection*{中文对应答案}
\StudyTarget{Compact___25}
\input{\StudyRoot/CS5489/exam-review/Midterm/items/MT005/answer-zh.tex}
\input{\StudyRoot/CS5489/exam-review/Midterm/items/MT005/concise-extra.tex}
\renewcommand{\StudyQuestionExtra}{\StudyTarget{Compact__mt034}\StudyTarget{Compact__l2q15}}
\StudyQuestion{MT034}{无限数据能消除分类错误吗}{简答题 / Written response · 判断并说明}
\input{\StudyRoot/CS5489/exam-review/Midterm/items/MT034/provenance.tex}
\subsection*{English question}
\StudyTarget{Compact__english-question_12}
\input{\StudyRoot/CS5489/exam-review/Midterm/items/MT034/question-en.tex}
\subsection*{中文题意}
\StudyTarget{Compact___26}
\input{\StudyRoot/CS5489/exam-review/Midterm/items/MT034/question-zh.tex}
\subsection*{English answer}
\StudyTarget{Compact__english-answer_12}
\input{\StudyRoot/CS5489/exam-review/Midterm/items/MT034/answer-en.tex}
\subsection*{中文对应答案}
\StudyTarget{Compact___27}
\input{\StudyRoot/CS5489/exam-review/Midterm/items/MT034/answer-zh.tex}
\input{\StudyRoot/CS5489/exam-review/Midterm/items/MT034/concise-extra.tex}
\renewcommand{\StudyQuestionExtra}{\StudyTarget{Compact__mt014}\StudyTarget{Compact__l2q03}}
\StudyQuestion{MT014}{生成式与判别式的区别}{选择题（多选） / Multiple select · 选不正确项 / Select incorrect statements}
\input{\StudyRoot/CS5489/exam-review/Midterm/items/MT014/provenance.tex}
\subsection*{English question}
\StudyTarget{Compact__english-question_13}
\input{\StudyRoot/CS5489/exam-review/Midterm/items/MT014/question-en.tex}
\subsection*{中文题意}
\StudyTarget{Compact___28}
\input{\StudyRoot/CS5489/exam-review/Midterm/items/MT014/question-zh.tex}
\subsection*{English answer}
\StudyTarget{Compact__english-answer_13}
\input{\StudyRoot/CS5489/exam-review/Midterm/items/MT014/answer-en.tex}
\subsection*{中文对应答案}
\StudyTarget{Compact___29}
\input{\StudyRoot/CS5489/exam-review/Midterm/items/MT014/answer-zh.tex}
\input{\StudyRoot/CS5489/exam-review/Midterm/items/MT014/concise-extra.tex}
\renewcommand{\StudyQuestionExtra}{\StudyTarget{Compact__mt033}\StudyTarget{Compact__l2q04}}
\StudyQuestion{MT033}{用图像特征解释两种建模方式}{简答题 / Written response · 比较、解释}
\input{\StudyRoot/CS5489/exam-review/Midterm/items/MT033/provenance.tex}
\subsection*{English question}
\StudyTarget{Compact__english-question_14}
\input{\StudyRoot/CS5489/exam-review/Midterm/items/MT033/question-en.tex}
\subsection*{中文题意}
\StudyTarget{Compact___30}
\input{\StudyRoot/CS5489/exam-review/Midterm/items/MT033/question-zh.tex}
\subsection*{English answer}
\StudyTarget{Compact__english-answer_14}
\input{\StudyRoot/CS5489/exam-review/Midterm/items/MT033/answer-en.tex}
\subsection*{中文对应答案}
\StudyTarget{Compact___31}
\input{\StudyRoot/CS5489/exam-review/Midterm/items/MT033/answer-zh.tex}
\input{\StudyRoot/CS5489/exam-review/Midterm/items/MT033/concise-extra.tex}
\renewcommand{\StudyQuestionExtra}{\StudyTarget{Compact__mt059}\StudyTarget{Compact__l2q13}\StudyTarget{Compact__mt059-lrgaussian-nb}}
\StudyQuestion{MT059}{LR与共享方差Gaussian NB}{简答题 / Written response · 比较}
\input{\StudyRoot/CS5489/exam-review/Midterm/items/MT059/provenance.tex}
\subsection*{English question}
\StudyTarget{Compact__english-question_15}
\input{\StudyRoot/CS5489/exam-review/Midterm/items/MT059/question-en.tex}
\subsection*{中文题意}
\StudyTarget{Compact___32}
\input{\StudyRoot/CS5489/exam-review/Midterm/items/MT059/question-zh.tex}
\subsection*{English answer}
\StudyTarget{Compact__english-answer_15}
\input{\StudyRoot/CS5489/exam-review/Midterm/items/MT059/answer-en.tex}
\subsection*{中文对应答案}
\StudyTarget{Compact___33}
\input{\StudyRoot/CS5489/exam-review/Midterm/items/MT059/answer-zh.tex}
\input{\StudyRoot/CS5489/exam-review/Midterm/items/MT059/concise-extra.tex}
\renewcommand{\StudyQuestionExtra}{\StudyTarget{Compact__mt021}\StudyTarget{Compact__l2q12}}
\StudyQuestion{MT021}{同为线性分类器，学习原则有何不同}{简答题 / Written response · 比较}
\input{\StudyRoot/CS5489/exam-review/Midterm/items/MT021/provenance.tex}
\subsection*{English question}
\StudyTarget{Compact__english-question_16}
\input{\StudyRoot/CS5489/exam-review/Midterm/items/MT021/question-en.tex}
\subsection*{中文题意}
\StudyTarget{Compact___34}
\input{\StudyRoot/CS5489/exam-review/Midterm/items/MT021/question-zh.tex}
\subsection*{English answer}
\StudyTarget{Compact__english-answer_16}
\input{\StudyRoot/CS5489/exam-review/Midterm/items/MT021/answer-en.tex}
\subsection*{中文对应答案}
\StudyTarget{Compact___35}
\input{\StudyRoot/CS5489/exam-review/Midterm/items/MT021/answer-zh.tex}
\input{\StudyRoot/CS5489/exam-review/Midterm/items/MT021/concise-extra.tex}
\renewcommand{\StudyQuestionExtra}{\StudyTarget{Compact__mt017}\StudyTarget{Compact__l2q11}}
\StudyQuestion{MT017}{哪些模型能产生不同形状的边界}{选择题（多选） / Multiple select}
\input{\StudyRoot/CS5489/exam-review/Midterm/items/MT017/provenance.tex}
\subsection*{English question}
\StudyTarget{Compact__english-question_17}
\input{\StudyRoot/CS5489/exam-review/Midterm/items/MT017/question-en.tex}
\subsection*{中文题意}
\StudyTarget{Compact___36}
\input{\StudyRoot/CS5489/exam-review/Midterm/items/MT017/question-zh.tex}
\subsection*{English answer}
\StudyTarget{Compact__english-answer_17}
\input{\StudyRoot/CS5489/exam-review/Midterm/items/MT017/answer-en.tex}
\subsection*{中文对应答案}
\StudyTarget{Compact___37}
\input{\StudyRoot/CS5489/exam-review/Midterm/items/MT017/answer-zh.tex}
\input{\StudyRoot/CS5489/exam-review/Midterm/items/MT017/concise-extra.tex}
\renewcommand{\StudyQuestionExtra}{\StudyTarget{Compact__mt027}\StudyTarget{Compact__l2q10}\StudyTarget{Compact__mt027-nb}}
\StudyQuestion{MT027}{连续特征能否使用NB：混合选择题}{选择题（多选） / Multiple select}
\input{\StudyRoot/CS5489/exam-review/Midterm/items/MT027/provenance.tex}
\subsection*{English question}
\StudyTarget{Compact__english-question_18}
\input{\StudyRoot/CS5489/exam-review/Midterm/items/MT027/question-en.tex}
\subsection*{中文题意}
\StudyTarget{Compact___38}
\input{\StudyRoot/CS5489/exam-review/Midterm/items/MT027/question-zh.tex}
\subsection*{English answer}
\StudyTarget{Compact__english-answer_18}
\input{\StudyRoot/CS5489/exam-review/Midterm/items/MT027/answer-en.tex}
\subsection*{中文对应答案}
\StudyTarget{Compact___39}
\input{\StudyRoot/CS5489/exam-review/Midterm/items/MT027/answer-zh.tex}
\input{\StudyRoot/CS5489/exam-review/Midterm/items/MT027/concise-extra.tex}
\renewcommand{\StudyQuestionExtra}{\StudyTarget{Compact__mt058}\StudyTarget{Compact__l2q16}\StudyTarget{Compact__mt058-nb}}
\StudyQuestion{MT058}{智能笔上的文本NB如何改进}{简答题 / Written response · 提出办法}
\input{\StudyRoot/CS5489/exam-review/Midterm/items/MT058/provenance.tex}
\subsection*{English question}
\StudyTarget{Compact__english-question_19}
\input{\StudyRoot/CS5489/exam-review/Midterm/items/MT058/question-en.tex}
\subsection*{中文题意}
\StudyTarget{Compact___40}
\input{\StudyRoot/CS5489/exam-review/Midterm/items/MT058/question-zh.tex}
\subsection*{English answer}
\StudyTarget{Compact__english-answer_19}
\input{\StudyRoot/CS5489/exam-review/Midterm/items/MT058/answer-en.tex}
\subsection*{中文对应答案}
\StudyTarget{Compact___41}
\input{\StudyRoot/CS5489/exam-review/Midterm/items/MT058/answer-zh.tex}
\input{\StudyRoot/CS5489/exam-review/Midterm/items/MT058/concise-extra.tex}
\Needspace{10\baselineskip}\phantomsection\addcontentsline{toc}{chapter}{逻辑回归与参数正则}
\renewcommand{\StudyQuestionExtra}{\StudyTarget{Compact__mt041}\StudyTarget{Compact__mt041-logistic}}
\StudyQuestion{MT041}{Logistic的概率、多分类与最优解}{选择题（多选） / Multiple select}
\input{\StudyRoot/CS5489/exam-review/Midterm/items/MT041/provenance.tex}
\subsection*{English question}
\StudyTarget{Compact__english-question_20}
\input{\StudyRoot/CS5489/exam-review/Midterm/items/MT041/question-en.tex}
\subsection*{中文题意}
\StudyTarget{Compact___42}
\input{\StudyRoot/CS5489/exam-review/Midterm/items/MT041/question-zh.tex}
\subsection*{English answer}
\StudyTarget{Compact__english-answer_20}
\input{\StudyRoot/CS5489/exam-review/Midterm/items/MT041/answer-en.tex}
\subsection*{中文对应答案}
\StudyTarget{Compact___43}
\input{\StudyRoot/CS5489/exam-review/Midterm/items/MT041/answer-zh.tex}
\input{\StudyRoot/CS5489/exam-review/Midterm/items/MT041/concise-extra.tex}
\renewcommand{\StudyQuestionExtra}{\StudyTarget{Compact__mt047}\StudyTarget{Compact__mt047-logistic}}
\StudyQuestion{MT047}{解释Logistic目标中的两项}{简答题 / Written response · 解释}
\input{\StudyRoot/CS5489/exam-review/Midterm/items/MT047/provenance.tex}
\subsection*{English question}
\StudyTarget{Compact__english-question_21}
\input{\StudyRoot/CS5489/exam-review/Midterm/items/MT047/question-en.tex}
\subsection*{中文题意}
\StudyTarget{Compact___44}
\input{\StudyRoot/CS5489/exam-review/Midterm/items/MT047/question-zh.tex}
\subsection*{English answer}
\StudyTarget{Compact__english-answer_21}
\input{\StudyRoot/CS5489/exam-review/Midterm/items/MT047/answer-en.tex}
\subsection*{中文对应答案}
\StudyTarget{Compact___45}
\input{\StudyRoot/CS5489/exam-review/Midterm/items/MT047/answer-zh.tex}
\input{\StudyRoot/CS5489/exam-review/Midterm/items/MT047/concise-extra.tex}
\renewcommand{\StudyQuestionExtra}{\StudyTarget{Compact__mt024}}
\StudyQuestion{MT024}{损失与正则项：\ensuremath{\alpha}从哪里来？}{简答题 / Written response · 解释}
\input{\StudyRoot/CS5489/exam-review/Midterm/items/MT024/provenance.tex}
\subsection*{English question}
\StudyTarget{Compact__english-question_22}
\input{\StudyRoot/CS5489/exam-review/Midterm/items/MT024/question-en.tex}
\subsection*{中文题意}
\StudyTarget{Compact___46}
\input{\StudyRoot/CS5489/exam-review/Midterm/items/MT024/question-zh.tex}
\subsection*{English answer}
\StudyTarget{Compact__english-answer_22}
\input{\StudyRoot/CS5489/exam-review/Midterm/items/MT024/answer-en.tex}
\subsection*{中文对应答案}
\StudyTarget{Compact___47}
\input{\StudyRoot/CS5489/exam-review/Midterm/items/MT024/answer-zh.tex}
\input{\StudyRoot/CS5489/exam-review/Midterm/items/MT024/concise-extra.tex}
\renewcommand{\StudyQuestionExtra}{\StudyTarget{Compact__mt018}\StudyTarget{Compact__mt018-l2}}
\StudyQuestion{MT018}{为什么在线性分类器上加L2？}{选择题（多选） / Multiple select}
\input{\StudyRoot/CS5489/exam-review/Midterm/items/MT018/provenance.tex}
\subsection*{English question}
\StudyTarget{Compact__english-question_23}
\input{\StudyRoot/CS5489/exam-review/Midterm/items/MT018/question-en.tex}
\subsection*{中文题意}
\StudyTarget{Compact___48}
\input{\StudyRoot/CS5489/exam-review/Midterm/items/MT018/question-zh.tex}
\subsection*{English answer}
\StudyTarget{Compact__english-answer_23}
\input{\StudyRoot/CS5489/exam-review/Midterm/items/MT018/answer-en.tex}
\subsection*{中文对应答案}
\StudyTarget{Compact___49}
\input{\StudyRoot/CS5489/exam-review/Midterm/items/MT018/answer-zh.tex}
\input{\StudyRoot/CS5489/exam-review/Midterm/items/MT018/concise-extra.tex}
\renewcommand{\StudyQuestionExtra}{\StudyTarget{Compact__mt012}\StudyTarget{Compact__mt012-l1}}
\StudyQuestion{MT012}{L1逻辑回归：写目标并画边界}{简答题 / Written response · 列式、作图、解释}
\input{\StudyRoot/CS5489/exam-review/Midterm/items/MT012/provenance.tex}
\subsection*{English question}
\StudyTarget{Compact__english-question_24}
\input{\StudyRoot/CS5489/exam-review/Midterm/items/MT012/question-en.tex}
\subsection*{中文题意}
\StudyTarget{Compact___50}
\input{\StudyRoot/CS5489/exam-review/Midterm/items/MT012/question-zh.tex}
\subsection*{English answer}
\StudyTarget{Compact__english-answer_24}
\input{\StudyRoot/CS5489/exam-review/Midterm/items/MT012/answer-en.tex}
\subsection*{中文对应答案}
\StudyTarget{Compact___51}
\input{\StudyRoot/CS5489/exam-review/Midterm/items/MT012/answer-zh.tex}
\input{\StudyRoot/CS5489/exam-review/Midterm/items/MT012/concise-extra.tex}
\renewcommand{\StudyQuestionExtra}{\StudyTarget{Compact__mt001}\StudyTarget{Compact__mt001-logistic}}
\StudyQuestion{MT001}{Logistic训练：哪些说法不成立？}{选择题（多选） / Multiple select · 选不正确项 / Select incorrect statements}
\input{\StudyRoot/CS5489/exam-review/Midterm/items/MT001/provenance.tex}
\subsection*{English question}
\StudyTarget{Compact__english-question_25}
\input{\StudyRoot/CS5489/exam-review/Midterm/items/MT001/question-en.tex}
\subsection*{中文题意}
\StudyTarget{Compact___52}
\input{\StudyRoot/CS5489/exam-review/Midterm/items/MT001/question-zh.tex}
\subsection*{English answer}
\StudyTarget{Compact__english-answer_25}
\input{\StudyRoot/CS5489/exam-review/Midterm/items/MT001/answer-en.tex}
\subsection*{中文对应答案}
\StudyTarget{Compact___53}
\input{\StudyRoot/CS5489/exam-review/Midterm/items/MT001/answer-zh.tex}
\input{\StudyRoot/CS5489/exam-review/Midterm/items/MT001/concise-extra.tex}
\renewcommand{\StudyQuestionExtra}{\StudyTarget{Compact__mt053}\StudyTarget{Compact__l2q14}\StudyTarget{Compact__mt053-lr}}
\StudyQuestion{MT053}{LR与生成式模型：混合选择题}{选择题（多选） / Multiple select}
\input{\StudyRoot/CS5489/exam-review/Midterm/items/MT053/provenance.tex}
\subsection*{English question}
\StudyTarget{Compact__english-question_26}
\input{\StudyRoot/CS5489/exam-review/Midterm/items/MT053/question-en.tex}
\subsection*{中文题意}
\StudyTarget{Compact___54}
\input{\StudyRoot/CS5489/exam-review/Midterm/items/MT053/question-zh.tex}
\subsection*{English answer}
\StudyTarget{Compact__english-answer_26}
\input{\StudyRoot/CS5489/exam-review/Midterm/items/MT053/answer-en.tex}
\subsection*{中文对应答案}
\StudyTarget{Compact___55}
\input{\StudyRoot/CS5489/exam-review/Midterm/items/MT053/answer-zh.tex}
\input{\StudyRoot/CS5489/exam-review/Midterm/items/MT053/concise-extra.tex}
\Needspace{10\baselineskip}\phantomsection\addcontentsline{toc}{chapter}{SVM与核方法}
\renewcommand{\StudyQuestionExtra}{\StudyTarget{Compact__mt042}\StudyTarget{Compact__mt042-svm}}
\StudyQuestion{MT042}{SVM的基本特点}{选择题（多选） / Multiple select}
\input{\StudyRoot/CS5489/exam-review/Midterm/items/MT042/provenance.tex}
\subsection*{English question}
\StudyTarget{Compact__english-question_27}
\input{\StudyRoot/CS5489/exam-review/Midterm/items/MT042/question-en.tex}
\subsection*{中文题意}
\StudyTarget{Compact___56}
\input{\StudyRoot/CS5489/exam-review/Midterm/items/MT042/question-zh.tex}
\subsection*{English answer}
\StudyTarget{Compact__english-answer_27}
\input{\StudyRoot/CS5489/exam-review/Midterm/items/MT042/answer-en.tex}
\subsection*{中文对应答案}
\StudyTarget{Compact___57}
\input{\StudyRoot/CS5489/exam-review/Midterm/items/MT042/answer-zh.tex}
\input{\StudyRoot/CS5489/exam-review/Midterm/items/MT042/concise-extra.tex}
\renewcommand{\StudyQuestionExtra}{\StudyTarget{Compact__mt048}\StudyTarget{Compact__mt048-kernel-trick}}
\StudyQuestion{MT048}{用一个例子讲清Kernel Trick}{简答题 / Written response · 解释、举例}
\input{\StudyRoot/CS5489/exam-review/Midterm/items/MT048/provenance.tex}
\subsection*{English question}
\StudyTarget{Compact__english-question_28}
\input{\StudyRoot/CS5489/exam-review/Midterm/items/MT048/question-en.tex}
\subsection*{中文题意}
\StudyTarget{Compact___58}
\input{\StudyRoot/CS5489/exam-review/Midterm/items/MT048/question-zh.tex}
\subsection*{English answer}
\StudyTarget{Compact__english-answer_28}
\input{\StudyRoot/CS5489/exam-review/Midterm/items/MT048/answer-en.tex}
\subsection*{中文对应答案}
\StudyTarget{Compact___59}
\input{\StudyRoot/CS5489/exam-review/Midterm/items/MT048/answer-zh.tex}
\input{\StudyRoot/CS5489/exam-review/Midterm/items/MT048/concise-extra.tex}
\renewcommand{\StudyQuestionExtra}{\StudyTarget{Compact__mt016}\StudyTarget{Compact__mt016-svm}}
\StudyQuestion{MT016}{核SVM：变换、存储与非向量输入}{选择题（多选） / Multiple select}
\input{\StudyRoot/CS5489/exam-review/Midterm/items/MT016/provenance.tex}
\subsection*{English question}
\StudyTarget{Compact__english-question_29}
\input{\StudyRoot/CS5489/exam-review/Midterm/items/MT016/question-en.tex}
\subsection*{中文题意}
\StudyTarget{Compact___60}
\input{\StudyRoot/CS5489/exam-review/Midterm/items/MT016/question-zh.tex}
\subsection*{English answer}
\StudyTarget{Compact__english-answer_29}
\input{\StudyRoot/CS5489/exam-review/Midterm/items/MT016/answer-en.tex}
\subsection*{中文对应答案}
\StudyTarget{Compact___61}
\input{\StudyRoot/CS5489/exam-review/Midterm/items/MT016/answer-zh.tex}
\input{\StudyRoot/CS5489/exam-review/Midterm/items/MT016/concise-extra.tex}
\renewcommand{\StudyQuestionExtra}{\StudyTarget{Compact__mt031}\StudyTarget{Compact__mt031-rbf}}
\StudyQuestion{MT031}{RBF距离与带宽：看近点和远点}{选择题（多选） / Multiple select}
\input{\StudyRoot/CS5489/exam-review/Midterm/items/MT031/provenance.tex}
\subsection*{English question}
\StudyTarget{Compact__english-question_30}
\input{\StudyRoot/CS5489/exam-review/Midterm/items/MT031/question-en.tex}
\subsection*{中文题意}
\StudyTarget{Compact___62}
\input{\StudyRoot/CS5489/exam-review/Midterm/items/MT031/question-zh.tex}
\subsection*{English answer}
\StudyTarget{Compact__english-answer_30}
\input{\StudyRoot/CS5489/exam-review/Midterm/items/MT031/answer-en.tex}
\subsection*{中文对应答案}
\StudyTarget{Compact___63}
\input{\StudyRoot/CS5489/exam-review/Midterm/items/MT031/answer-zh.tex}
\input{\StudyRoot/CS5489/exam-review/Midterm/items/MT031/concise-extra.tex}
\renewcommand{\StudyQuestionExtra}{\StudyTarget{Compact__mt055}}
\StudyQuestion{MT055}{判断五个候选核是否合法}{选择题（多选） / Multiple select}
\input{\StudyRoot/CS5489/exam-review/Midterm/items/MT055/provenance.tex}
\subsection*{English question}
\StudyTarget{Compact__english-question_31}
\input{\StudyRoot/CS5489/exam-review/Midterm/items/MT055/question-en.tex}
\subsection*{中文题意}
\StudyTarget{Compact___64}
\input{\StudyRoot/CS5489/exam-review/Midterm/items/MT055/question-zh.tex}
\subsection*{English answer}
\StudyTarget{Compact__english-answer_31}
\input{\StudyRoot/CS5489/exam-review/Midterm/items/MT055/answer-en.tex}
\subsection*{中文对应答案}
\StudyTarget{Compact___65}
\input{\StudyRoot/CS5489/exam-review/Midterm/items/MT055/answer-zh.tex}
\input{\StudyRoot/CS5489/exam-review/Midterm/items/MT055/concise-extra.tex}
\renewcommand{\StudyQuestionExtra}{\StudyTarget{Compact__mt060}\StudyTarget{Compact__mt060-svmcc}}
\StudyQuestion{MT060}{二次核SVM：大C与小C的边界}{简答题 / Written response · 作图、比较、解释}
\input{\StudyRoot/CS5489/exam-review/Midterm/items/MT060/provenance.tex}
\subsection*{English question}
\StudyTarget{Compact__english-question_32}
\input{\StudyRoot/CS5489/exam-review/Midterm/items/MT060/question-en.tex}
\subsection*{中文题意}
\StudyTarget{Compact___66}
\input{\StudyRoot/CS5489/exam-review/Midterm/items/MT060/question-zh.tex}
\subsection*{English answer}
\StudyTarget{Compact__english-answer_32}
\input{\StudyRoot/CS5489/exam-review/Midterm/items/MT060/answer-en.tex}
\subsection*{中文对应答案}
\StudyTarget{Compact___67}
\input{\StudyRoot/CS5489/exam-review/Midterm/items/MT060/answer-zh.tex}
\input{\StudyRoot/CS5489/exam-review/Midterm/items/MT060/concise-extra.tex}
\renewcommand{\StudyQuestionExtra}{\StudyTarget{Compact__mt054}\StudyTarget{Compact__mt054-crbf}}
\StudyQuestion{MT054}{C、支持向量与RBF的零训练误差}{选择题（多选） / Multiple select}
\input{\StudyRoot/CS5489/exam-review/Midterm/items/MT054/provenance.tex}
\subsection*{English question}
\StudyTarget{Compact__english-question_33}
\input{\StudyRoot/CS5489/exam-review/Midterm/items/MT054/question-en.tex}
\subsection*{中文题意}
\StudyTarget{Compact___68}
\input{\StudyRoot/CS5489/exam-review/Midterm/items/MT054/question-zh.tex}
\subsection*{English answer}
\StudyTarget{Compact__english-answer_33}
\input{\StudyRoot/CS5489/exam-review/Midterm/items/MT054/answer-en.tex}
\subsection*{中文对应答案}
\StudyTarget{Compact___69}
\input{\StudyRoot/CS5489/exam-review/Midterm/items/MT054/answer-zh.tex}
\input{\StudyRoot/CS5489/exam-review/Midterm/items/MT054/concise-extra.tex}
\renewcommand{\StudyQuestionExtra}{\StudyTarget{Compact__mt043}\StudyTarget{Compact__mt043-svm}}
\StudyQuestion{MT043}{核SVM怎样处理非线性？}{选择题（多选） / Multiple select}
\input{\StudyRoot/CS5489/exam-review/Midterm/items/MT043/provenance.tex}
\subsection*{English question}
\StudyTarget{Compact__english-question_34}
\input{\StudyRoot/CS5489/exam-review/Midterm/items/MT043/question-en.tex}
\subsection*{中文题意}
\StudyTarget{Compact___70}
\input{\StudyRoot/CS5489/exam-review/Midterm/items/MT043/question-zh.tex}
\subsection*{English answer}
\StudyTarget{Compact__english-answer_34}
\input{\StudyRoot/CS5489/exam-review/Midterm/items/MT043/answer-en.tex}
\subsection*{中文对应答案}
\StudyTarget{Compact___71}
\input{\StudyRoot/CS5489/exam-review/Midterm/items/MT043/answer-zh.tex}
\input{\StudyRoot/CS5489/exam-review/Midterm/items/MT043/concise-extra.tex}
\renewcommand{\StudyQuestionExtra}{\StudyTarget{Compact__mt029}\StudyTarget{Compact__mt029-svm}}
\StudyQuestion{MT029}{哪些SVM说法错误？}{选择题（多选） / Multiple select · 选不正确项 / Select incorrect statements}
\input{\StudyRoot/CS5489/exam-review/Midterm/items/MT029/provenance.tex}
\subsection*{English question}
\StudyTarget{Compact__english-question_35}
\input{\StudyRoot/CS5489/exam-review/Midterm/items/MT029/question-en.tex}
\subsection*{中文题意}
\StudyTarget{Compact___72}
\input{\StudyRoot/CS5489/exam-review/Midterm/items/MT029/question-zh.tex}
\subsection*{English answer}
\StudyTarget{Compact__english-answer_35}
\input{\StudyRoot/CS5489/exam-review/Midterm/items/MT029/answer-en.tex}
\subsection*{中文对应答案}
\StudyTarget{Compact___73}
\input{\StudyRoot/CS5489/exam-review/Midterm/items/MT029/answer-zh.tex}
\input{\StudyRoot/CS5489/exam-review/Midterm/items/MT029/concise-extra.tex}
\renewcommand{\StudyQuestionExtra}{\StudyTarget{Compact__mt010}\StudyTarget{Compact__mt010-gaussian-bayessvm}}
\StudyQuestion{MT010}{同一张分布图，Gaussian Bayes与二次核SVM如何判断？}{简答题 / Written response · 读图或公式、比较}
\input{\StudyRoot/CS5489/exam-review/Midterm/items/MT010/provenance.tex}
\subsection*{English question}
\StudyTarget{Compact__english-question_36}
\input{\StudyRoot/CS5489/exam-review/Midterm/items/MT010/question-en.tex}
\subsection*{中文题意}
\StudyTarget{Compact___74}
\input{\StudyRoot/CS5489/exam-review/Midterm/items/MT010/question-zh.tex}
\subsection*{English answer}
\StudyTarget{Compact__english-answer_36}
\input{\StudyRoot/CS5489/exam-review/Midterm/items/MT010/answer-en.tex}
\subsection*{中文对应答案}
\StudyTarget{Compact___75}
\input{\StudyRoot/CS5489/exam-review/Midterm/items/MT010/answer-zh.tex}
\input{\StudyRoot/CS5489/exam-review/Midterm/items/MT010/concise-extra.tex}
\renewcommand{\StudyQuestionExtra}{\StudyTarget{Compact__mt009}\StudyTarget{Compact__mt009-svm}}
\StudyQuestion{MT009}{十万维特征、两千样本：SVM怎样提速？}{简答题 / Written response · 提出办法、解释}
\input{\StudyRoot/CS5489/exam-review/Midterm/items/MT009/provenance.tex}
\subsection*{English question}
\StudyTarget{Compact__english-question_37}
\input{\StudyRoot/CS5489/exam-review/Midterm/items/MT009/question-en.tex}
\subsection*{中文题意}
\StudyTarget{Compact___76}
\input{\StudyRoot/CS5489/exam-review/Midterm/items/MT009/question-zh.tex}
\subsection*{English answer}
\StudyTarget{Compact__english-answer_37}
\input{\StudyRoot/CS5489/exam-review/Midterm/items/MT009/answer-en.tex}
\subsection*{中文对应答案}
\StudyTarget{Compact___77}
\input{\StudyRoot/CS5489/exam-review/Midterm/items/MT009/answer-zh.tex}
\input{\StudyRoot/CS5489/exam-review/Midterm/items/MT009/concise-extra.tex}
\renewcommand{\StudyQuestionExtra}{\StudyTarget{Compact__mt061}\StudyTarget{Compact__mt061-svmrbf-svm}}
\StudyQuestion{MT061}{预测内存：线性SVM与RBF SVM}{简答题 / Written response · 排序、解释}
\input{\StudyRoot/CS5489/exam-review/Midterm/items/MT061/provenance.tex}
\subsection*{English question}
\StudyTarget{Compact__english-question_38}
\input{\StudyRoot/CS5489/exam-review/Midterm/items/MT061/question-en.tex}
\subsection*{中文题意}
\StudyTarget{Compact___78}
\input{\StudyRoot/CS5489/exam-review/Midterm/items/MT061/question-zh.tex}
\subsection*{English answer}
\StudyTarget{Compact__english-answer_38}
\input{\StudyRoot/CS5489/exam-review/Midterm/items/MT061/answer-en.tex}
\subsection*{中文对应答案}
\StudyTarget{Compact___79}
\input{\StudyRoot/CS5489/exam-review/Midterm/items/MT061/answer-zh.tex}
\input{\StudyRoot/CS5489/exam-review/Midterm/items/MT061/concise-extra.tex}
\Needspace{10\baselineskip}\phantomsection\addcontentsline{toc}{chapter}{回归、特征选择与残差损失}
\renewcommand{\StudyQuestionExtra}{\StudyTarget{Compact__mt019}\StudyTarget{Compact__mt019-ridgelasso}}
\StudyQuestion{MT019}{Ridge、LASSO与零系数}{选择题（多选） / Multiple select}
\input{\StudyRoot/CS5489/exam-review/Midterm/items/MT019/provenance.tex}
\subsection*{English question}
\StudyTarget{Compact__english-question_39}
\input{\StudyRoot/CS5489/exam-review/Midterm/items/MT019/question-en.tex}
\subsection*{中文题意}
\StudyTarget{Compact___80}
\input{\StudyRoot/CS5489/exam-review/Midterm/items/MT019/question-zh.tex}
\subsection*{English answer}
\StudyTarget{Compact__english-answer_39}
\input{\StudyRoot/CS5489/exam-review/Midterm/items/MT019/answer-en.tex}
\subsection*{中文对应答案}
\StudyTarget{Compact___81}
\input{\StudyRoot/CS5489/exam-review/Midterm/items/MT019/answer-zh.tex}
\input{\StudyRoot/CS5489/exam-review/Midterm/items/MT019/concise-extra.tex}
\renewcommand{\StudyQuestionExtra}{\StudyTarget{Compact__mt045}\StudyTarget{Compact__mt045-ridgeomp}}
\StudyQuestion{MT045}{特征选择：零系数、Ridge与OMP}{选择题（多选） / Multiple select}
\input{\StudyRoot/CS5489/exam-review/Midterm/items/MT045/provenance.tex}
\subsection*{English question}
\StudyTarget{Compact__english-question_40}
\input{\StudyRoot/CS5489/exam-review/Midterm/items/MT045/question-en.tex}
\subsection*{中文题意}
\StudyTarget{Compact___82}
\input{\StudyRoot/CS5489/exam-review/Midterm/items/MT045/question-zh.tex}
\subsection*{English answer}
\StudyTarget{Compact__english-answer_40}
\input{\StudyRoot/CS5489/exam-review/Midterm/items/MT045/answer-en.tex}
\subsection*{中文对应答案}
\StudyTarget{Compact___83}
\input{\StudyRoot/CS5489/exam-review/Midterm/items/MT045/answer-zh.tex}
\input{\StudyRoot/CS5489/exam-review/Midterm/items/MT045/concise-extra.tex}
\renewcommand{\StudyQuestionExtra}{\StudyTarget{Compact__mt030}\StudyTarget{Compact__mt030-l1l2}}
\StudyQuestion{MT030}{L1/L2回归的益处与措辞边界}{选择题（多选） / Multiple select}
\input{\StudyRoot/CS5489/exam-review/Midterm/items/MT030/provenance.tex}
\subsection*{English question}
\StudyTarget{Compact__english-question_41}
\input{\StudyRoot/CS5489/exam-review/Midterm/items/MT030/question-en.tex}
\subsection*{中文题意}
\StudyTarget{Compact___84}
\input{\StudyRoot/CS5489/exam-review/Midterm/items/MT030/question-zh.tex}
\subsection*{English answer}
\StudyTarget{Compact__english-answer_41}
\input{\StudyRoot/CS5489/exam-review/Midterm/items/MT030/answer-en.tex}
\subsection*{中文对应答案}
\StudyTarget{Compact___85}
\input{\StudyRoot/CS5489/exam-review/Midterm/items/MT030/answer-zh.tex}
\input{\StudyRoot/CS5489/exam-review/Midterm/items/MT030/concise-extra.tex}
\renewcommand{\StudyQuestionExtra}{\StudyTarget{Compact__mt057}\StudyTarget{Compact__mt057-l1l2elastic-net}}
\StudyQuestion{MT057}{同时使用L1与L2：Elastic Net}{选择题（多选） / Multiple select}
\input{\StudyRoot/CS5489/exam-review/Midterm/items/MT057/provenance.tex}
\subsection*{English question}
\StudyTarget{Compact__english-question_42}
\input{\StudyRoot/CS5489/exam-review/Midterm/items/MT057/question-en.tex}
\subsection*{中文题意}
\StudyTarget{Compact___86}
\input{\StudyRoot/CS5489/exam-review/Midterm/items/MT057/question-zh.tex}
\subsection*{English answer}
\StudyTarget{Compact__english-answer_42}
\input{\StudyRoot/CS5489/exam-review/Midterm/items/MT057/answer-en.tex}
\subsection*{中文对应答案}
\StudyTarget{Compact___87}
\input{\StudyRoot/CS5489/exam-review/Midterm/items/MT057/answer-zh.tex}
\input{\StudyRoot/CS5489/exam-review/Midterm/items/MT057/concise-extra.tex}
\renewcommand{\StudyQuestionExtra}{\StudyTarget{Compact__mt025}}
\StudyQuestion{MT025}{绝对残差损失对异常值有什么影响？}{简答题 / Written response · 解释}
\input{\StudyRoot/CS5489/exam-review/Midterm/items/MT025/provenance.tex}
\subsection*{English question}
\StudyTarget{Compact__english-question_43}
\input{\StudyRoot/CS5489/exam-review/Midterm/items/MT025/question-en.tex}
\subsection*{中文题意}
\StudyTarget{Compact___88}
\input{\StudyRoot/CS5489/exam-review/Midterm/items/MT025/question-zh.tex}
\subsection*{English answer}
\StudyTarget{Compact__english-answer_43}
\input{\StudyRoot/CS5489/exam-review/Midterm/items/MT025/answer-en.tex}
\subsection*{中文对应答案}
\StudyTarget{Compact___89}
\input{\StudyRoot/CS5489/exam-review/Midterm/items/MT025/answer-zh.tex}
\input{\StudyRoot/CS5489/exam-review/Midterm/items/MT025/concise-extra.tex}
\renewcommand{\StudyQuestionExtra}{\StudyTarget{Compact__mt063}\StudyTarget{Compact__mt063-huber}}
\StudyQuestion{MT063}{Huber损失：画图并解释好处}{简答题 / Written response · 作图、解释}
\input{\StudyRoot/CS5489/exam-review/Midterm/items/MT063/provenance.tex}
\subsection*{English question}
\StudyTarget{Compact__english-question_44}
\input{\StudyRoot/CS5489/exam-review/Midterm/items/MT063/question-en.tex}
\subsection*{中文题意}
\StudyTarget{Compact___90}
\input{\StudyRoot/CS5489/exam-review/Midterm/items/MT063/question-zh.tex}
\subsection*{English answer}
\StudyTarget{Compact__english-answer_44}
\input{\StudyRoot/CS5489/exam-review/Midterm/items/MT063/answer-en.tex}
\subsection*{中文对应答案}
\StudyTarget{Compact___91}
\input{\StudyRoot/CS5489/exam-review/Midterm/items/MT063/answer-zh.tex}
\input{\StudyRoot/CS5489/exam-review/Midterm/items/MT063/concise-extra.tex}
\renewcommand{\StudyQuestionExtra}{\StudyTarget{Compact__mt013}\StudyTarget{Compact__mt013-l1l2}}
\StudyQuestion{MT013}{残差L1＋权重L2，各自在做什么？}{简答题 / Written response · 解释}
\input{\StudyRoot/CS5489/exam-review/Midterm/items/MT013/provenance.tex}
\subsection*{English question}
\StudyTarget{Compact__english-question_45}
\input{\StudyRoot/CS5489/exam-review/Midterm/items/MT013/question-en.tex}
\subsection*{中文题意}
\StudyTarget{Compact___92}
\input{\StudyRoot/CS5489/exam-review/Midterm/items/MT013/question-zh.tex}
\subsection*{English answer}
\StudyTarget{Compact__english-answer_45}
\input{\StudyRoot/CS5489/exam-review/Midterm/items/MT013/answer-en.tex}
\subsection*{中文对应答案}
\StudyTarget{Compact___93}
\input{\StudyRoot/CS5489/exam-review/Midterm/items/MT013/answer-zh.tex}
\input{\StudyRoot/CS5489/exam-review/Midterm/items/MT013/concise-extra.tex}
\renewcommand{\StudyQuestionExtra}{\StudyTarget{Compact__mt050}}
\StudyQuestion{MT050}{缺车比车多更糟：设计非对称损失}{简答题 / Written response · 设计、作图、解释}
\input{\StudyRoot/CS5489/exam-review/Midterm/items/MT050/provenance.tex}
\subsection*{English question}
\StudyTarget{Compact__english-question_46}
\input{\StudyRoot/CS5489/exam-review/Midterm/items/MT050/question-en.tex}
\subsection*{中文题意}
\StudyTarget{Compact___94}
\input{\StudyRoot/CS5489/exam-review/Midterm/items/MT050/question-zh.tex}
\subsection*{English answer}
\StudyTarget{Compact__english-answer_46}
\input{\StudyRoot/CS5489/exam-review/Midterm/items/MT050/answer-en.tex}
\subsection*{中文对应答案}
\StudyTarget{Compact___95}
\input{\StudyRoot/CS5489/exam-review/Midterm/items/MT050/answer-zh.tex}
\input{\StudyRoot/CS5489/exam-review/Midterm/items/MT050/concise-extra.tex}
\renewcommand{\StudyQuestionExtra}{\StudyTarget{Compact__mt037}}
\StudyQuestion{MT037}{线性回归为何与均值预测器表现相近？}{简答题 / Written response · 提出办法、解释}
\input{\StudyRoot/CS5489/exam-review/Midterm/items/MT037/provenance.tex}
\subsection*{English question}
\StudyTarget{Compact__english-question_47}
\input{\StudyRoot/CS5489/exam-review/Midterm/items/MT037/question-en.tex}
\subsection*{中文题意}
\StudyTarget{Compact___96}
\input{\StudyRoot/CS5489/exam-review/Midterm/items/MT037/question-zh.tex}
\subsection*{English answer}
\StudyTarget{Compact__english-answer_47}
\input{\StudyRoot/CS5489/exam-review/Midterm/items/MT037/answer-en.tex}
\subsection*{中文对应答案}
\StudyTarget{Compact___97}
\input{\StudyRoot/CS5489/exam-review/Midterm/items/MT037/answer-zh.tex}
\input{\StudyRoot/CS5489/exam-review/Midterm/items/MT037/concise-extra.tex}
\Needspace{10\baselineskip}\phantomsection\addcontentsline{toc}{chapter}{集成学习}
\renewcommand{\StudyQuestionExtra}{\StudyTarget{Compact__mt006}\StudyTarget{Compact__mt006-baggingboosting}}
\StudyQuestion{MT006}{Bagging与Boosting的基本区别}{选择题（多选） / Multiple select}
\input{\StudyRoot/CS5489/exam-review/Midterm/items/MT006/provenance.tex}
\subsection*{English question}
\StudyTarget{Compact__english-question_48}
\input{\StudyRoot/CS5489/exam-review/Midterm/items/MT006/question-en.tex}
\subsection*{中文题意}
\StudyTarget{Compact___98}
\input{\StudyRoot/CS5489/exam-review/Midterm/items/MT006/question-zh.tex}
\subsection*{English answer}
\StudyTarget{Compact__english-answer_48}
\input{\StudyRoot/CS5489/exam-review/Midterm/items/MT006/answer-en.tex}
\subsection*{中文对应答案}
\StudyTarget{Compact___99}
\input{\StudyRoot/CS5489/exam-review/Midterm/items/MT006/answer-zh.tex}
\input{\StudyRoot/CS5489/exam-review/Midterm/items/MT006/concise-extra.tex}
\renewcommand{\StudyQuestionExtra}{\StudyTarget{Compact__mt044}\StudyTarget{Compact__mt044-boosting}}
\StudyQuestion{MT044}{随机森林与Boosting的用途}{选择题（多选） / Multiple select}
\input{\StudyRoot/CS5489/exam-review/Midterm/items/MT044/provenance.tex}
\subsection*{English question}
\StudyTarget{Compact__english-question_49}
\input{\StudyRoot/CS5489/exam-review/Midterm/items/MT044/question-en.tex}
\subsection*{中文题意}
\StudyTarget{Compact___100}
\input{\StudyRoot/CS5489/exam-review/Midterm/items/MT044/question-zh.tex}
\subsection*{English answer}
\StudyTarget{Compact__english-answer_49}
\input{\StudyRoot/CS5489/exam-review/Midterm/items/MT044/answer-en.tex}
\subsection*{中文对应答案}
\StudyTarget{Compact___101}
\input{\StudyRoot/CS5489/exam-review/Midterm/items/MT044/answer-zh.tex}
\input{\StudyRoot/CS5489/exam-review/Midterm/items/MT044/concise-extra.tex}
\renewcommand{\StudyQuestionExtra}{\StudyTarget{Compact__mt062}}
\StudyQuestion{MT062}{随机森林：树多了为什么仍有共同误差？}{简答题 / Written response · 读图或公式、提出办法}
\input{\StudyRoot/CS5489/exam-review/Midterm/items/MT062/provenance.tex}
\subsection*{English question}
\StudyTarget{Compact__english-question_50}
\input{\StudyRoot/CS5489/exam-review/Midterm/items/MT062/question-en.tex}
\subsection*{中文题意}
\StudyTarget{Compact___102}
\input{\StudyRoot/CS5489/exam-review/Midterm/items/MT062/question-zh.tex}
\subsection*{English answer}
\StudyTarget{Compact__english-answer_50}
\input{\StudyRoot/CS5489/exam-review/Midterm/items/MT062/answer-en.tex}
\subsection*{中文对应答案}
\StudyTarget{Compact___103}
\input{\StudyRoot/CS5489/exam-review/Midterm/items/MT062/answer-zh.tex}
\input{\StudyRoot/CS5489/exam-review/Midterm/items/MT062/concise-extra.tex}
\renewcommand{\StudyQuestionExtra}{\StudyTarget{Compact__mt056}}
\StudyQuestion{MT056}{集成方法：树深、并行与后续轮次}{选择题（多选） / Multiple select}
\input{\StudyRoot/CS5489/exam-review/Midterm/items/MT056/provenance.tex}
\subsection*{English question}
\StudyTarget{Compact__english-question_51}
\input{\StudyRoot/CS5489/exam-review/Midterm/items/MT056/question-en.tex}
\subsection*{中文题意}
\StudyTarget{Compact___104}
\input{\StudyRoot/CS5489/exam-review/Midterm/items/MT056/question-zh.tex}
\subsection*{English answer}
\StudyTarget{Compact__english-answer_51}
\input{\StudyRoot/CS5489/exam-review/Midterm/items/MT056/answer-en.tex}
\subsection*{中文对应答案}
\StudyTarget{Compact___105}
\input{\StudyRoot/CS5489/exam-review/Midterm/items/MT056/answer-zh.tex}
\input{\StudyRoot/CS5489/exam-review/Midterm/items/MT056/concise-extra.tex}
\renewcommand{\StudyQuestionExtra}{\StudyTarget{Compact__mt038}\StudyTarget{Compact__mt038-100rfadaboost}}
\StudyQuestion{MT038}{100核部署：RF与AdaBoost如何并行预测？}{简答题 / Written response · 比较、提出办法}
\input{\StudyRoot/CS5489/exam-review/Midterm/items/MT038/provenance.tex}
\subsection*{English question}
\StudyTarget{Compact__english-question_52}
\input{\StudyRoot/CS5489/exam-review/Midterm/items/MT038/question-en.tex}
\subsection*{中文题意}
\StudyTarget{Compact___106}
\input{\StudyRoot/CS5489/exam-review/Midterm/items/MT038/question-zh.tex}
\subsection*{English answer}
\StudyTarget{Compact__english-answer_52}
\input{\StudyRoot/CS5489/exam-review/Midterm/items/MT038/answer-en.tex}
\subsection*{中文对应答案}
\StudyTarget{Compact___107}
\input{\StudyRoot/CS5489/exam-review/Midterm/items/MT038/answer-zh.tex}
\input{\StudyRoot/CS5489/exam-review/Midterm/items/MT038/concise-extra.tex}
\renewcommand{\StudyQuestionExtra}{\StudyTarget{Compact__mt049}\StudyTarget{Compact__mt049-adaboost}}
\StudyQuestion{MT049}{在线更新AdaBoost：速度与异常值}{简答题 / Written response · 提出办法、判断并说明}
\input{\StudyRoot/CS5489/exam-review/Midterm/items/MT049/provenance.tex}
\subsection*{English question}
\StudyTarget{Compact__english-question_53}
\input{\StudyRoot/CS5489/exam-review/Midterm/items/MT049/question-en.tex}
\subsection*{中文题意}
\StudyTarget{Compact___108}
\input{\StudyRoot/CS5489/exam-review/Midterm/items/MT049/question-zh.tex}
\subsection*{English answer}
\StudyTarget{Compact__english-answer_53}
\input{\StudyRoot/CS5489/exam-review/Midterm/items/MT049/answer-en.tex}
\subsection*{中文对应答案}
\StudyTarget{Compact___109}
\input{\StudyRoot/CS5489/exam-review/Midterm/items/MT049/answer-zh.tex}
\input{\StudyRoot/CS5489/exam-review/Midterm/items/MT049/concise-extra.tex}
\renewcommand{\StudyQuestionExtra}{\StudyTarget{Compact__mt064}\StudyTarget{Compact__mt064-baggingboostinghuber}}
\StudyQuestion{MT064}{综合题：Bagging/Boosting与Huber}{简答题 / Written response · 比较、作图、解释}
\input{\StudyRoot/CS5489/exam-review/Midterm/items/MT064/provenance.tex}
\subsection*{English question}
\StudyTarget{Compact__english-question_54}
\input{\StudyRoot/CS5489/exam-review/Midterm/items/MT064/question-en.tex}
\subsection*{中文题意}
\StudyTarget{Compact___110}
\input{\StudyRoot/CS5489/exam-review/Midterm/items/MT064/question-zh.tex}
\subsection*{English answer}
\StudyTarget{Compact__english-answer_54}
\input{\StudyRoot/CS5489/exam-review/Midterm/items/MT064/answer-en.tex}
\subsection*{中文对应答案}
\StudyTarget{Compact___111}
\input{\StudyRoot/CS5489/exam-review/Midterm/items/MT064/answer-zh.tex}
\input{\StudyRoot/CS5489/exam-review/Midterm/items/MT064/concise-extra.tex}
\Needspace{10\baselineskip}\phantomsection\addcontentsline{toc}{chapter}{分类损失图与错误代价}
\renewcommand{\StudyQuestionExtra}{\StudyTarget{Compact__mt008}}
\StudyQuestion{MT008}{平方形分类损失为什么惩罚“太正确”？}{简答题 / Written response · 读图或公式、判断并说明}
\input{\StudyRoot/CS5489/exam-review/Midterm/items/MT008/provenance.tex}
\subsection*{English question}
\StudyTarget{Compact__english-question_55}
\input{\StudyRoot/CS5489/exam-review/Midterm/items/MT008/question-en.tex}
\subsection*{中文题意}
\StudyTarget{Compact___112}
\input{\StudyRoot/CS5489/exam-review/Midterm/items/MT008/question-zh.tex}
\subsection*{English answer}
\StudyTarget{Compact__english-answer_55}
\input{\StudyRoot/CS5489/exam-review/Midterm/items/MT008/answer-en.tex}
\subsection*{中文对应答案}
\StudyTarget{Compact___113}
\input{\StudyRoot/CS5489/exam-review/Midterm/items/MT008/answer-zh.tex}
\input{\StudyRoot/CS5489/exam-review/Midterm/items/MT008/concise-extra.tex}
\renewcommand{\StudyQuestionExtra}{\StudyTarget{Compact__mt026}\StudyTarget{Compact__mt026-vz1}}
\StudyQuestion{MT026}{V形损失：最优点为何在z=1？}{简答题 / Written response · 读图或公式、判断并说明}
\input{\StudyRoot/CS5489/exam-review/Midterm/items/MT026/provenance.tex}
\subsection*{English question}
\StudyTarget{Compact__english-question_56}
\input{\StudyRoot/CS5489/exam-review/Midterm/items/MT026/question-en.tex}
\subsection*{中文题意}
\StudyTarget{Compact___114}
\input{\StudyRoot/CS5489/exam-review/Midterm/items/MT026/question-zh.tex}
\subsection*{English answer}
\StudyTarget{Compact__english-answer_56}
\input{\StudyRoot/CS5489/exam-review/Midterm/items/MT026/answer-en.tex}
\subsection*{中文对应答案}
\StudyTarget{Compact___115}
\input{\StudyRoot/CS5489/exam-review/Midterm/items/MT026/answer-zh.tex}
\input{\StudyRoot/CS5489/exam-review/Midterm/items/MT026/concise-extra.tex}
\renewcommand{\StudyQuestionExtra}{\StudyTarget{Compact__mt051}}
\StudyQuestion{MT051}{负间隔也零损失：分类器会放过什么？}{简答题 / Written response · 读图或公式、判断并说明、作图}
\input{\StudyRoot/CS5489/exam-review/Midterm/items/MT051/provenance.tex}
\subsection*{English question}
\StudyTarget{Compact__english-question_57}
\input{\StudyRoot/CS5489/exam-review/Midterm/items/MT051/question-en.tex}
\subsection*{中文题意}
\StudyTarget{Compact___116}
\input{\StudyRoot/CS5489/exam-review/Midterm/items/MT051/question-zh.tex}
\subsection*{English answer}
\StudyTarget{Compact__english-answer_57}
\input{\StudyRoot/CS5489/exam-review/Midterm/items/MT051/answer-en.tex}
\subsection*{中文对应答案}
\StudyTarget{Compact___117}
\input{\StudyRoot/CS5489/exam-review/Midterm/items/MT051/answer-zh.tex}
\input{\StudyRoot/CS5489/exam-review/Midterm/items/MT051/concise-extra.tex}
\renewcommand{\StudyQuestionExtra}{\StudyTarget{Compact__mt039}}
\StudyQuestion{MT039}{饱和的非凸分类损失}{简答题 / Written response · 读图或公式}
\input{\StudyRoot/CS5489/exam-review/Midterm/items/MT039/provenance.tex}
\subsection*{English question}
\StudyTarget{Compact__english-question_58}
\input{\StudyRoot/CS5489/exam-review/Midterm/items/MT039/question-en.tex}
\subsection*{中文题意}
\StudyTarget{Compact___118}
\input{\StudyRoot/CS5489/exam-review/Midterm/items/MT039/question-zh.tex}
\subsection*{English answer}
\StudyTarget{Compact__english-answer_58}
\input{\StudyRoot/CS5489/exam-review/Midterm/items/MT039/answer-en.tex}
\subsection*{中文对应答案}
\StudyTarget{Compact___119}
\input{\StudyRoot/CS5489/exam-review/Midterm/items/MT039/answer-zh.tex}
\input{\StudyRoot/CS5489/exam-review/Midterm/items/MT039/concise-extra.tex}
\renewcommand{\StudyQuestionExtra}{\StudyTarget{Compact__mt065}\StudyTarget{Compact__mt065-logistichingeexponential}}
\StudyQuestion{MT065}{比较Logistic、Hinge和Exponential损失}{简答题 / Written response · 比较、排序}
\input{\StudyRoot/CS5489/exam-review/Midterm/items/MT065/provenance.tex}
\subsection*{English question}
\StudyTarget{Compact__english-question_59}
\input{\StudyRoot/CS5489/exam-review/Midterm/items/MT065/question-en.tex}
\subsection*{中文题意}
\StudyTarget{Compact___120}
\input{\StudyRoot/CS5489/exam-review/Midterm/items/MT065/question-zh.tex}
\subsection*{English answer}
\StudyTarget{Compact__english-answer_59}
\input{\StudyRoot/CS5489/exam-review/Midterm/items/MT065/answer-en.tex}
\subsection*{中文对应答案}
\StudyTarget{Compact___121}
\input{\StudyRoot/CS5489/exam-review/Midterm/items/MT065/answer-zh.tex}
\input{\StudyRoot/CS5489/exam-review/Midterm/items/MT065/concise-extra.tex}
\renewcommand{\StudyQuestionExtra}{\StudyTarget{Compact__mt007}}
\StudyQuestion{MT007}{筛查任务：漏诊比误报更贵}{简答题 / Written response · 提出办法、解释}
\input{\StudyRoot/CS5489/exam-review/Midterm/items/MT007/provenance.tex}
\subsection*{English question}
\StudyTarget{Compact__english-question_60}
\input{\StudyRoot/CS5489/exam-review/Midterm/items/MT007/question-en.tex}
\subsection*{中文题意}
\StudyTarget{Compact___122}
\input{\StudyRoot/CS5489/exam-review/Midterm/items/MT007/question-zh.tex}
\subsection*{English answer}
\StudyTarget{Compact__english-answer_60}
\input{\StudyRoot/CS5489/exam-review/Midterm/items/MT007/answer-en.tex}
\subsection*{中文对应答案}
\StudyTarget{Compact___123}
\input{\StudyRoot/CS5489/exam-review/Midterm/items/MT007/answer-zh.tex}
\input{\StudyRoot/CS5489/exam-review/Midterm/items/MT007/concise-extra.tex}
\renewcommand{\StudyQuestionExtra}{\StudyTarget{Compact__mt023}\StudyTarget{Compact__mt023-100050}}
\StudyQuestion{MT023}{1000名受试者中只有50名阳性}{简答题 / Written response · 诊断问题、提出办法}
\input{\StudyRoot/CS5489/exam-review/Midterm/items/MT023/provenance.tex}
\subsection*{English question}
\StudyTarget{Compact__english-question_61}
\input{\StudyRoot/CS5489/exam-review/Midterm/items/MT023/question-en.tex}
\subsection*{中文题意}
\StudyTarget{Compact___124}
\input{\StudyRoot/CS5489/exam-review/Midterm/items/MT023/question-zh.tex}
\subsection*{English answer}
\StudyTarget{Compact__english-answer_61}
\input{\StudyRoot/CS5489/exam-review/Midterm/items/MT023/answer-en.tex}
\subsection*{中文对应答案}
\StudyTarget{Compact___125}
\input{\StudyRoot/CS5489/exam-review/Midterm/items/MT023/answer-zh.tex}
\input{\StudyRoot/CS5489/exam-review/Midterm/items/MT023/concise-extra.tex}
\renewcommand{\StudyQuestionExtra}{\StudyTarget{Compact__mt035}}
\StudyQuestion{MT035}{正常邮件被大量拦截怎么办？}{简答题 / Written response · 提出办法、解释}
\input{\StudyRoot/CS5489/exam-review/Midterm/items/MT035/provenance.tex}
\subsection*{English question}
\StudyTarget{Compact__english-question_62}
\input{\StudyRoot/CS5489/exam-review/Midterm/items/MT035/question-en.tex}
\subsection*{中文题意}
\StudyTarget{Compact___126}
\input{\StudyRoot/CS5489/exam-review/Midterm/items/MT035/question-zh.tex}
\subsection*{English answer}
\StudyTarget{Compact__english-answer_62}
\input{\StudyRoot/CS5489/exam-review/Midterm/items/MT035/answer-en.tex}
\subsection*{中文对应答案}
\StudyTarget{Compact___127}
\input{\StudyRoot/CS5489/exam-review/Midterm/items/MT035/answer-zh.tex}
\input{\StudyRoot/CS5489/exam-review/Midterm/items/MT035/concise-extra.tex}
\Needspace{10\baselineskip}\phantomsection\addcontentsline{toc}{chapter}{神经网络与链式法则}
\renewcommand{\StudyQuestionExtra}{\StudyTarget{Compact__mt067}}
\StudyQuestion{MT067}{链式求导与两层网络反传}{简答题 / Written response · 计算推导}
\input{\StudyRoot/CS5489/exam-review/Midterm/items/MT067/provenance.tex}
\subsection*{English question}
\StudyTarget{Compact__english-question_63}
\input{\StudyRoot/CS5489/exam-review/Midterm/items/MT067/question-en.tex}
\subsection*{中文题意}
\StudyTarget{Compact___128}
\input{\StudyRoot/CS5489/exam-review/Midterm/items/MT067/question-zh.tex}
\subsection*{English answer}
\StudyTarget{Compact__english-answer_63}
\input{\StudyRoot/CS5489/exam-review/Midterm/items/MT067/answer-en.tex}
\subsection*{中文对应答案}
\StudyTarget{Compact___129}
\input{\StudyRoot/CS5489/exam-review/Midterm/items/MT067/answer-zh.tex}
\input{\StudyRoot/CS5489/exam-review/Midterm/items/MT067/concise-extra.tex}
\renewcommand{\StudyQuestionExtra}{\StudyTarget{Compact__mt066}}
\StudyQuestion{MT066}{梯度消失：为什么深度和激活函数有关？}{简答题 / Written response · 解释、计算推导}
\input{\StudyRoot/CS5489/exam-review/Midterm/items/MT066/provenance.tex}
\subsection*{English question}
\StudyTarget{Compact__english-question_64}
\input{\StudyRoot/CS5489/exam-review/Midterm/items/MT066/question-en.tex}
\subsection*{中文题意}
\StudyTarget{Compact___130}
\input{\StudyRoot/CS5489/exam-review/Midterm/items/MT066/question-zh.tex}
\subsection*{English answer}
\StudyTarget{Compact__english-answer_64}
\input{\StudyRoot/CS5489/exam-review/Midterm/items/MT066/answer-en.tex}
\subsection*{中文对应答案}
\StudyTarget{Compact___131}
\input{\StudyRoot/CS5489/exam-review/Midterm/items/MT066/answer-zh.tex}
\input{\StudyRoot/CS5489/exam-review/Midterm/items/MT066/concise-extra.tex}
\Needspace{10\baselineskip}\phantomsection\addcontentsline{toc}{chapter}{历史范围拓展：Gaussian Process}
\renewcommand{\StudyQuestionExtra}{\StudyTarget{Compact__mt003}\StudyTarget{Compact__mt003-gaussian-process}}
\StudyQuestion{MT003}{历史拓展：Gaussian Process的假设}{选择题（多选） / Multiple select}
\input{\StudyRoot/CS5489/exam-review/Midterm/items/MT003/provenance.tex}
\subsection*{English question}
\StudyTarget{Compact__english-question_65}
\input{\StudyRoot/CS5489/exam-review/Midterm/items/MT003/question-en.tex}
\subsection*{中文题意}
\StudyTarget{Compact___132}
\input{\StudyRoot/CS5489/exam-review/Midterm/items/MT003/question-zh.tex}
\subsection*{English answer}
\StudyTarget{Compact__english-answer_65}
\input{\StudyRoot/CS5489/exam-review/Midterm/items/MT003/answer-en.tex}
\subsection*{中文对应答案}
\StudyTarget{Compact___133}
\input{\StudyRoot/CS5489/exam-review/Midterm/items/MT003/answer-zh.tex}
\input{\StudyRoot/CS5489/exam-review/Midterm/items/MT003/concise-extra.tex}
\renewcommand{\StudyQuestionExtra}{\StudyTarget{Compact__mt036}\StudyTarget{Compact__mt036-gprf}}
\StudyQuestion{MT036}{历史拓展：音乐地理回归，GP还是RF？}{简答题 / Written response · 比较}
\input{\StudyRoot/CS5489/exam-review/Midterm/items/MT036/provenance.tex}
\subsection*{English question}
\StudyTarget{Compact__english-question_66}
\input{\StudyRoot/CS5489/exam-review/Midterm/items/MT036/question-en.tex}
\subsection*{中文题意}
\StudyTarget{Compact___134}
\input{\StudyRoot/CS5489/exam-review/Midterm/items/MT036/question-zh.tex}
\subsection*{English answer}
\StudyTarget{Compact__english-answer_66}
\input{\StudyRoot/CS5489/exam-review/Midterm/items/MT036/answer-en.tex}
\subsection*{中文对应答案}
\StudyTarget{Compact___135}
\input{\StudyRoot/CS5489/exam-review/Midterm/items/MT036/answer-zh.tex}
\input{\StudyRoot/CS5489/exam-review/Midterm/items/MT036/concise-extra.tex}
\clearpage\StudyHeading{1}{原卷索引 / Original-paper index}
按原卷题号查找本册位置；不同年份的重复题指向同一MT编号。2021B*保留题纸与配套答案的年份冲突。
\StudyTarget{Compact__original-paper-index}
\subsection*{2020B期中Quiz}
\StudyTarget{Compact__2020bquiz}
\begin{longtable}{@{}p{.2\linewidth}p{.56\linewidth}p{.18\linewidth}@{}}
\toprule 原题 & 本册题号与页码 & 分值 \\ \midrule\endhead
Q1 & \StudyRef{Compact}{mt001}{MT001} & 5 \\
Q2 & \StudyRef{Compact}{mt002}{MT002} & 5 \\
Q3 & \StudyRef{Compact}{mt003}{MT003} & 5 \\
Q4 & \StudyRef{Compact}{mt004}{MT004} & 5 \\
Q5 & \StudyRef{Compact}{mt005}{MT005} & 5 \\
Q6 & \StudyRef{Compact}{mt006}{MT006} & 5 \\
Q7 & \StudyRef{Compact}{mt007}{MT007} & 10 \\
Q8 & \StudyRef{Compact}{mt008}{MT008} & 10 \\
Q9 & \StudyRef{Compact}{mt009}{MT009} & 10 \\
Q10 & \StudyRef{Compact}{mt010}{MT010} & 10 \\
Q11 & \StudyRef{Compact}{mt011}{MT011} & 10 \\
Q12 & \StudyRef{Compact}{mt012}{MT012} & 10 \\
Q13 & \StudyRef{Compact}{mt013}{MT013} & 10 \\
\bottomrule\end{longtable}
\subsection*{2021A期中}
\StudyTarget{Compact__2021a}
\begin{longtable}{@{}p{.2\linewidth}p{.56\linewidth}p{.18\linewidth}@{}}
\toprule 原题 & 本册题号与页码 & 分值 \\ \midrule\endhead
Q1 & \StudyRef{Compact}{mt014}{MT014} & 5 \\
Q2 & \StudyRef{Compact}{mt015}{MT015} & 5 \\
Q3 & \StudyRef{Compact}{mt016}{MT016} & 5 \\
Q4 & \StudyRef{Compact}{mt017}{MT017} & 5 \\
Q5 & \StudyRef{Compact}{mt018}{MT018} & 5 \\
Q6 & \StudyRef{Compact}{mt019}{MT019} & 5 \\
Q7 & \StudyRef{Compact}{mt020}{MT020} & 10 \\
Q8 & \StudyRef{Compact}{mt021}{MT021} & 10 \\
Q9 & \StudyRef{Compact}{mt022}{MT022} & 10 \\
Q10 & \StudyRef{Compact}{mt023}{MT023} & 10 \\
Q11 & \StudyRef{Compact}{mt024}{MT024} & 10 \\
Q12 & \StudyRef{Compact}{mt025}{MT025} & 10 \\
Q13 & \StudyRef{Compact}{mt026}{MT026} & 10 \\
\bottomrule\end{longtable}
\subsection*{2021B*期中}
\StudyTarget{Compact__2021b}
\begin{longtable}{@{}p{.2\linewidth}p{.56\linewidth}p{.18\linewidth}@{}}
\toprule 原题 & 本册题号与页码 & 分值 \\ \midrule\endhead
Q1 & \StudyRef{Compact}{mt027}{MT027} & 5 \\
Q2 & \StudyRef{Compact}{mt028}{MT028} & 5 \\
Q3 & \StudyRef{Compact}{mt029}{MT029} & 5 \\
Q4 & \StudyRef{Compact}{mt030}{MT030} & 5 \\
Q5 & \StudyRef{Compact}{mt031}{MT031} & 5 \\
Q6 & \StudyRef{Compact}{mt032}{MT032} & 5 \\
Q7 & \StudyRef{Compact}{mt033}{MT033} & 10 \\
Q8 & \StudyRef{Compact}{mt034}{MT034} & 10 \\
Q9 & \StudyRef{Compact}{mt035}{MT035} & 10 \\
Q10 & \StudyRef{Compact}{mt036}{MT036} & 10 \\
Q11 & \StudyRef{Compact}{mt037}{MT037} & 10 \\
Q12 & \StudyRef{Compact}{mt038}{MT038} & 10 \\
Q13 & \StudyRef{Compact}{mt039}{MT039} & 10 \\
\bottomrule\end{longtable}
\subsection*{2023A期中}
\StudyTarget{Compact__2023a}
\begin{longtable}{@{}p{.2\linewidth}p{.56\linewidth}p{.18\linewidth}@{}}
\toprule 原题 & 本册题号与页码 & 分值 \\ \midrule\endhead
Q1 & \StudyRef{Compact}{mt040}{MT040} & 5 \\
Q2 & \StudyRef{Compact}{mt041}{MT041} & 5 \\
Q3 & \StudyRef{Compact}{mt042}{MT042} & 5 \\
Q4 & \StudyRef{Compact}{mt043}{MT043} & 5 \\
Q5 & \StudyRef{Compact}{mt044}{MT044} & 5 \\
Q6 & \StudyRef{Compact}{mt045}{MT045} & 5 \\
Q7 & \StudyRef{Compact}{mt046}{MT046} & 10 \\
Q8 & \StudyRef{Compact}{mt047}{MT047} & 10 \\
Q9 & \StudyRef{Compact}{mt048}{MT048} & 10 \\
Q10 & \StudyRef{Compact}{mt023}{MT023} & 10 \\
Q11 & \StudyRef{Compact}{mt049}{MT049} & 10 \\
Q12 & \StudyRef{Compact}{mt050}{MT050} & 10 \\
Q13 & \StudyRef{Compact}{mt051}{MT051} & 10 \\
\bottomrule\end{longtable}
\subsection*{2023B期中}
\StudyTarget{Compact__2023b}
\begin{longtable}{@{}p{.2\linewidth}p{.56\linewidth}p{.18\linewidth}@{}}
\toprule 原题 & 本册题号与页码 & 分值 \\ \midrule\endhead
Q1 & \StudyRef{Compact}{mt052}{MT052} & 5 \\
Q2 & \StudyRef{Compact}{mt053}{MT053} & 5 \\
Q3 & \StudyRef{Compact}{mt054}{MT054} & 5 \\
Q4 & \StudyRef{Compact}{mt055}{MT055} & 5 \\
Q5 & \StudyRef{Compact}{mt056}{MT056} & 5 \\
Q6 & \StudyRef{Compact}{mt057}{MT057} & 5 \\
Q7 & \StudyRef{Compact}{mt058}{MT058} & 10 \\
Q8 & \StudyRef{Compact}{mt059}{MT059} & 10 \\
Q9 & \StudyRef{Compact}{mt060}{MT060} & 10 \\
Q10 & \StudyRef{Compact}{mt061}{MT061} & 10 \\
Q11 & \StudyRef{Compact}{mt062}{MT062} & 10 \\
Q12 & \StudyRef{Compact}{mt063}{MT063} & 10 \\
Q13 & \StudyRef{Compact}{mt007}{MT007} & 10 \\
\bottomrule\end{longtable}
\subsection*{2025A期中}
\StudyTarget{Compact__2025a}
\begin{longtable}{@{}p{.2\linewidth}p{.56\linewidth}p{.18\linewidth}@{}}
\toprule 原题 & 本册题号与页码 & 分值 \\ \midrule\endhead
Q1 & \StudyRef{Compact}{mt040}{MT040} & 5 \\
Q2 & \StudyRef{Compact}{mt041}{MT041} & 5 \\
Q3 & \StudyRef{Compact}{mt042}{MT042} & 5 \\
Q4 & \StudyRef{Compact}{mt043}{MT043} & 5 \\
Q5 & \StudyRef{Compact}{mt044}{MT044} & 5 \\
Q6 & \StudyRef{Compact}{mt057}{MT057} & 5 \\
Q7 & \StudyRef{Compact}{mt064}{MT064} & 10 \\
Q8 & \StudyRef{Compact}{mt047}{MT047} & 10 \\
Q9 & \StudyRef{Compact}{mt065}{MT065} & 10 \\
Q10 & \StudyRef{Compact}{mt066}{MT066} & 10 \\
Q11 & \StudyRef{Compact}{mt046}{MT046} & 10 \\
Q12 & \StudyRef{Compact}{mt050}{MT050} & 10 \\
Q13 & \StudyRef{Compact}{mt067}{MT067} & 10 \\
\bottomrule\end{longtable}
\end{document}
